
\documentclass[10pt]{article}
\usepackage[utf8]{inputenc}
\usepackage[a4paper, left=2.54cm, right=2.54cm, top=3cm, bottom=3cm]{geometry} % Márgenes personalizados a 2cm
\usepackage[spanish]{babel} % Para fechas y formato en español
\usepackage{multicol}
\usepackage{enumitem}
\usepackage{fancyhdr}
\usepackage{titlesec}

% Configuración de página
\pagestyle{fancy}
\fancyhf{} % Limpia encabezados y pies de página
\fancyhead[L]{Academia Prepolicial - MCD y MCM}
\fancyhead[R]{\today}
\fancyfoot[C]{\thepage} % Número de página centrado en el pie
\renewcommand{\headrulewidth}{0.8pt}

% Configuración de títulos y espaciado
\titleformat{\section}[block]{\large\bfseries\centering}{\thesection}{1em}{}
\setlength{\columnsep}{1cm}
\setlength{\parskip}{0.8em}
\setlength{\parindent}{0em}


\begin{document}

	\begin{center}
		\Large\bfseries Máximo Común Divisor y Mínimo Común Múltiplo
	\end{center}

\begin{multicols}{2}
	\section*{Ejercicios}
		\begin{enumerate}[label=\textbf{\arabic*.}, itemsep=0.4cm]
			\item Calcula el MCD y MCM de \( A \), \( B \) y \( C \) y da como respuesta la suma de ambos resultados. \\
			\( A = 3600 \) \\
			\( B = 180 \) \\
			\( C = 500 \)

			\item Calcula el MCD y MCM de \( A \), \( B \) y \( C \) y da como respuesta la suma de los exponentes de ambos resultados. \\
			\( A = 2^3 \times 3^5 \times 7^2 \) \\
			\( B = 2^4 \times 3^4 \times 5^3 \) \\
			\( C = 2^5 \times 5^4 \times 11^5 \)

			\item Si el MCM de \( A \) y \( B \) tiene 81 divisores, calcula \( n + 1 \): \\
			\( A = 12^n \times 3 \) \\
			\( B = 3^n \times 48 \)

			\item ¿Cuántos divisores tiene el MCD (\( A \), \( B \))? \\
			\( A = 48 \times 32 \times 73 \) \\
			\( B = 83 \times 27 \times 49 \)

			\item ¿Cuántos divisores tiene el MCD (\( A \), \( B \))? \\
			\( A = 4^6 \times 3^3 \times 7^4 \) \\
			\( B = 8^3 \times 16 \times 343 \)

			\item Calcula \( k \): \\
			\( \text{MCD}(210k, 300k, 420k) = 1200 \)

			\item Calcula \( n \): \\
			\( A = 6 \times 14^n \) \\
			\( B = 6^n \times 14 \) \\
			\( \text{MCD}(A, B) = 336 \)

			\item Determina el valor de \( K \): \\
			\( \text{MCM}\left( \frac{21k}{5}, \frac{7k}{10}, \frac{9k}{5} \right) = 1260 \)

			\item Calcula \( k + 1 \): \\
			\( \text{MCM}\left( \frac{13k}{7}, \frac{5k}{14}, \frac{8k}{7} \right) = 520 \)

			\item En un corral hay cierto número de gallinas que no pasan de 350 ni bajan de 380. Si las gallinas se acomodan en grupos de 2, 3, 4 o 5, siempre sobra 1. ¿Cuántas gallinas hay en el corral si se añaden 6?

			\item Tres aviones (\( A \), \( B \) y \( C \)) parten de una base a las 8 horas. Si \( A \) regresa cada hora y cuarto; \( B \), cada \( \frac{3}{4} \) de hora; y \( C \), cada 50 minutos, ¿a qué hora se reencontrarán por primera vez en la base?

			\item Calcula \( A + B \) si el MCD de ambos números es 495 y el valor de \( B \) se encuentra entre 5000 y 6000: \\
			\( A = a48b \) \\
			\( B = mnnm \)

			\item Calcula \( (a + b) \): \\
			\( \text{MCD}(a1b8, a9b0) = \degree^{88} \)

			\item Calcula la cifra de las unidades del mayor número de 3 cifras, que convertido a los sistemas de numeración de bases 6, 8 y 9 da como resultados números que terminan en 5, 7 y 8 respectivamente.




			% Para añadir más ejercicios, simplemente copia y pega el siguiente bloque:
			% \item [Enunciado del ejercicio]
		\end{enumerate}

\end{multicols}

\end{document}